% ============================================================
% EPISTEMIC STANDARD (COMPLETE FINAL VERSION)
% ============================================================
% Version: 1.0 Final
% Date: 2026-01-30
% Purpose: Complete epistemic framework for EDC Book 2
% Integration: Chapter 0 or Preface
% ============================================================

\section*{On Certainty, Verification, and Epistemic Status}
\addcontentsline{toc}{section}{On Certainty, Verification, and Epistemic Status}

% ------------------------------------------------------------
% What We Know With High Confidence
% ------------------------------------------------------------

\subsection*{What We Know With High Confidence}

Within its stated postulates [P], the EDC framework \textbf{yields} specific quantitative results.
Where a result is mapped to a standard observable, the corresponding model output
\textbf{matches current measurements} at the stated accuracy.

\textbf{Critical caveat (dictionary dependence):}
Any statement of the form
\[
X_{\mathrm{EDC}} \equiv Y_{\mathrm{BL}}
\]
requires a reduction dictionary (5D $\to$ 3D identification), unit/normalization conventions,
and boundary-condition choices. Such statements are therefore \textbf{conditional}.

% ------------------------------------------------------------
% Primary Sources of Uncertainty
% ------------------------------------------------------------

\subsection*{Primary Sources of Uncertainty}

Absolute certainty is not claimed because several components remain incomplete or under active development:

\begin{enumerate}
\item \textbf{Brane reduction physics is not fully closed}
\begin{itemize}
\item Boundary conditions are not uniquely fixed in the current presentation
\item Normalization conventions admit nontrivial choices
\item A reduction normalization factor $C_{\mathrm{red}}$ may enter (currently unconstrained)
\end{itemize}

\item \textbf{Higher-order effects are not yet computed for all observables}
\begin{itemize}
\item Electromagnetic and loop-level corrections (e.g.\ $O(\alpha)$)
\item Finite-size and compositeness effects where applicable
\item Renormalization-group running and scheme dependence (where relevant)
\end{itemize}

\item \textbf{Fundamental postulates are hypotheses}
\begin{itemize}
\item 5D membrane structure and reduction mechanism [P]
\item Discrete symmetry/crystallography assumptions (e.g.\ $Z_6$) [P]
\item Identification of specific topological defects with observed sectors [P]
\end{itemize}
\end{enumerate}

% ------------------------------------------------------------
% Mathematical Precision Levels (Subtags)
% ------------------------------------------------------------

\subsection*{Mathematical Precision Levels (Subtags)}

To avoid ambiguity in phrases like ``mathematically exact'', we use \textbf{precision subtags}
that refine (but do not replace) the core epistemic tags.

\begin{itemize}
\item \textbf{[Der:Sym] --- Exact Symbolic Derivation:}
closed-form algebraic expressions within the formal system (no truncation).

\item \textbf{[Der:Num] --- Exact Numerical Procedure:}
a fully specified deterministic numerical computation with a stated tolerance (e.g.\ $10^{-10}$),
validated by \textbf{three independent stability checks}:
\begin{enumerate}[(i)]
\item grid refinement to convergence (Richardson extrapolation),
\item cross-check via independent integrator (e.g.\ Gauss--Kronrod vs Simpson), and
\item asymptotic/analytic limit verification where applicable.
\end{enumerate}

\item \textbf{[Dc:Approx] --- Conditional, Exact Within an Explicit Approximation:}
results that are exact at a stated order/limit (e.g.\ tree level, leading order, thin-brane limit),
with missing corrections tracked as [Open].
\end{itemize}

% ------------------------------------------------------------
% Epistemic Tag Definitions (Formal)
% ------------------------------------------------------------

\subsection*{Epistemic Tag Definitions (Formal)}

\subsubsection*{[M] --- Mathematics}
Pure mathematical statement independent of physics.
\begin{itemize}
\item Criterion: provable within a formal system
\item Status: absolutely certain
\end{itemize}

\subsubsection*{[P] --- Postulate}
Physical hypothesis about the 5D structure and mechanism.
\begin{itemize}
\item Criterion: assumed starting point, not derived
\item Status: testable by consequences
\end{itemize}

\subsubsection*{[Der] --- Pure Derivation}
A mathematical consequence within the EDC formal system using only [M] and [P],
\textbf{without} interpreting the result as a baseline observable.
\begin{itemize}
\item Key rule: no claim of the form ``$X_{\mathrm{EDC}} = Y_{\mathrm{BL}}$''
\item Status: logically necessary \emph{given} the postulates
\end{itemize}

\subsubsection*{[Dc] --- Conditional Derivation (Default for Physics Outputs)}
A claim that connects EDC formal quantities to a baseline observable.
\begin{itemize}
\item Criterion: requires a reduction dictionary (5D $\to$ 3D), normalization, units, and conventions
\item Key boundary:
\begin{center}
\fbox{\parbox{0.92\textwidth}{
\textbf{Rule:} ANY claim of the form ``$X_{\mathrm{EDC}} = Y_{\mathrm{BL}}$'' is tagged [Dc],
because the identification invokes physical interpretation, normalization/units,
and reduction assumptions that are not purely mathematical.
}}
\end{center}
\item Status: conditional on dictionary and reduction physics (including $C_{\mathrm{red}}$ if applicable)
\end{itemize}

\subsubsection*{[I] --- Identified}
Structural/formal similarity recognition \textbf{without} numerical parameter tuning.
\begin{itemize}
\item Key boundary: if a parameter is adjusted to match data, it is not [I]
\end{itemize}

\subsubsection*{[Cal] --- Calibrated}
Explicit numerical adjustment of one or more free parameters to fit a baseline value.
\begin{itemize}
\item Example: ``we set parameter $p$ to match $Y_{\mathrm{BL}}$''
\end{itemize}

\subsubsection*{[BL] --- Baseline}
Empirical measurement or consensus value used as reference (e.g.\ CODATA/PDG/NASA tables).

\subsubsection*{[Open] --- Open Problem}
Unresolved issue, missing correction, or incomplete derivation step that materially affects certainty.

% ------------------------------------------------------------
% Meaning of [Dc] in Practice
% ------------------------------------------------------------

\subsection*{Meaning of [Dc] in Practice}

When a result is tagged [Dc], it should be read as follows:

\begin{itemize}
\item \textbf{NOT:} ``proved with absolute certainty''
\item \textbf{BUT:} ``derived within the model and mapped to an observable, with agreement
against the stated baseline, conditional on explicit reduction assumptions''
\end{itemize}

% ------------------------------------------------------------
% Agreement, Error Budgets, and Sensitivity
% ------------------------------------------------------------

\subsection*{Agreement, Error Budgets, and Sensitivity}

For each major EDC--BL comparison, we track an explicit correction envelope.
A minimal template is:

\begin{itemize}
\item \textbf{Baseline reference:} $Y_{\mathrm{BL}}$ is taken from the baseline constants table 
(Table~\ref{tab:baseline_constants}, PDG 2024 / CODATA 2022).

\item \textbf{Reduction normalization:} an explicit factor $C_{\mathrm{red}}$ is carried where 
relevant (currently [Open] if not derived from 5D action).

\item \textbf{Error budget:} leading missing corrections are listed with estimated magnitudes 
(e.g.\ EM $O(\alpha) \sim 0.1\%$, RG running $\sim 0.01\%$); unknowns remain [Open]. 
Total expected correction envelope is typically stated as a range $[A, B]\%$, against which 
the observed difference is compared.

\item \textbf{Sensitivity:} ``explored variations'' means parameters $p_i$ were varied over 
stated ranges
\[
p_i \in [p_i^{(0)}(1-\epsilon_i),\, p_i^{(0)}(1+\epsilon_i)]
\]
(e.g.\ $\sigma: \epsilon = 10\%$, $\ell_p/r_e: \epsilon = 1\%$), and the induced shift 
was recorded. Dominant uncertainty is typically $C_{\mathrm{red}}$ [Open].
\end{itemize}

% ------------------------------------------------------------
% Scientific Posture
% ------------------------------------------------------------

\subsection*{Scientific Posture}

This work adopts an explicit epistemic framework:
\begin{itemize}
\item Strong agreement claims are made where warranted by the stated model+dictionary
\item Incompleteness is bounded and stated, with missing pieces tracked as [Open]
\item Falsification and ``tension'' criteria are formulated relative to the correction envelope
\item Claims are tagged to distinguish axioms, derivations, identifications, and calibrations
\end{itemize}

This approach prioritizes \textbf{transparency over certainty claims}.

% ------------------------------------------------------------
% Falsification and Tension Criteria
% ------------------------------------------------------------

\subsection*{Falsification and Tension Criteria}

\textbf{Definitions:}

\begin{itemize}
\item \textbf{Tension:} Measured value lies outside estimated correction envelope 
  (error budget not fully closed; [Open] terms remain).

\item \textbf{Falsification:} After error budget closure (all [Open] corrections calculated), 
  residual deviation exceeds $3\sigma$ measurement uncertainty and no plausible 
  $C_{\mathrm{red}}$ restores agreement, \textit{or} fundamental geometric assumptions 
  proven incompatible with independent tests.
\end{itemize}

\textbf{Current status:}
\begin{itemize}
\item No results in tension (all predictions within expected correction envelopes)
\item Framework not falsified (no geometric contradictions found)
\item Validated pending error budget closure
\end{itemize}

\textbf{Path to definitive test:}
Once error budget closed (all [Open] corrections calculated):
\begin{itemize}
\item If agreement persists $\to$ Framework strongly confirmed
\item If tension emerges $\to$ Indicates missing physics or incorrect $C_{\mathrm{red}}$
\item If falsified $\to$ Geometric postulates require revision
\end{itemize}

% ============================================================
% End of Epistemic Standard Section
% ============================================================
